\documentclass{article}
\usepackage[margin=2cm]{geometry}
\usepackage{amsmath}
\usepackage{graphicx}
\usepackage{booktabs}
\usepackage[utf8]{inputenc}
\begin{document}
\subsection*{0.1 BPS Equations}
We now use the vanishing of the fermionic variations to obtain BPS equations for the warp factor and the three non-zero scalars.
\subsection*{0.1.1 Dilatino equation and projector}
We begin by imposing the vanishing of the dilatino variation, \(\delta \chi_{A}=0\), which implies
\[
\frac{1}{2} \gamma^{5} \sigma^{\prime} \varepsilon_{A}=N_{0} \varepsilon_{A}+N_{3} \gamma^{7} \left(\sigma^{3}\right)^{B} \quad_{A} \varepsilon_{B}
\]
This equation can be interpreted as a projection condition on the spinors \(\varepsilon_{A}\). Consistency of this projection condition then requires that
\[
\sigma^{\prime}=2 \eta \sqrt{N_{0}^{2}+N_{3}^{2}}
\]
where \(\eta= \pm 1\). Plugging this BPS equation back into then yields a second form of the projection condition,
\[
\gamma^{\tilde{5}} \varepsilon_{A}=G_{0} \varepsilon_{A}-G_{3} \gamma^{7}\left(\sigma^{3}\right)^{B} \varepsilon_{B}
\]
which is more useful in the derivation of the other BPS equations. In the above, we have defined
\[
G_{0}=\eta \frac{N_{0}}{\sqrt{N_{0}^{2}+N_{3}^{2}}} \quad \quad G_{3}=-\eta \frac{N_{3}}{\sqrt{N_{0}^{2}+N_{3}^{2}}}
\]
\subsection*{0.1.2 Gravitino equation}
The analysis of the gravitino equation \(\delta \psi_{A \mu}=0\) proceeds in exactly the same way as for the Lorentzian case studied in. The procedure gives rise to a first-order equation for the warp factor \(f\) and an algebraic constraint. To avoid excessive overlap with that paper, we simply cite the result,
\[
f^{\prime}=2\left(G_{0} S_{0}+G_{3} S_{3}\right) \quad e^{-2 f}=4\left(G_{0} S_{0}+G_{3} S_{3}\rangle^{2}-4\left(S_{0}^{2}+S_{3}^{2}\right)\right.
\]
\subsection*{0.1.3 Gaugino equations}
Finally, we turn toward the gaugino equation \(\delta \lambda_{A}^{\prime}=0\). Again the analysis of this equation proceeds in an exactly analogous manner to the Lorentzian case. The result is
\[
\cos \phi^{3}\left(\phi^{0}\right)^{\prime}=-\left(G_{0} M_{0}+G_{3} M_{3}\right) \quad\left(\phi^{3}\right)^{\prime}=i\left(G_{3} M_{0}-G_{0} M_{3}\right)
\]
The right-hand sides of both equations are real, and thus give rise to real solutions when appropriate initial conditions are imposed.
\end{document}
